\section{Conclusion and Future Work}
\label{sec:conclusion}

In this work, we have addressed the fundamental challenge of low-data calibration overfitting in dynamic weight-space model merging. By rejecting ad-hoc heuristic regularizers in favor of a mathematically rigorous approach, we derived the first formal Rademacher complexity generalization bound for dynamically merged models. Our learning-theory analysis proved that a model's generalization capability is directly scaled by the parameter-space distances (Frobenius or Spectral norms) of its task-specific experts from the pre-trained base model.

Guided by these theoretical derivations, we introduced Spectral and Rademacher-guided Routing Regularization (SR3), a novel, asymmetric regularization framework that scales the weight decay of routing parameters proportionally to their corresponding task-vector norms, as well as a direct, smoothed $L_1$ Group-Lasso formulation (SR3-L1). On a continuous weight-merging simulator featuring structured task-vector geometries and representation entanglement, we demonstrate that non-parametric methods collapse completely, while our proposed Frobenius (SR3-F), Spectral (SR3-S), and smoothed $L_1$ Group-Lasso (SR3-L1) variants achieve highly competitive joint multi-task accuracy comparable to state-of-the-art heuristics, while providing formal first-principles generalization guarantees. Under our simulated environment with structured task-vector geometries, we show that the spectral variant SR3-S indeed outperforms the Frobenius variant SR3-F, validating that constraining the worst-case representation distortion is a tighter generalization constraint than constraining average distortion. 

\textbf{Limitations and Future Work.}
While our findings are theoretically elegant and empirically robust, several avenues remain for future exploration:
\begin{enumerate}
    \item \textbf{Immediate Priority: Evaluating on Physical LLMs and Giant Foundation Models:} While our continuous simulator models a deep backbone, our immediate priority and immediate next step is to evaluate SR3 on physical large language models (LLMs, such as LLaMA-3 or Mistral) and vision-language models with billions of parameters under realistic multi-task settings. Under modern PEFT/LoRA adapter merging, task experts are fine-tuned as low-rank matrices. In these settings, computing exact task-vector spectral norms via standard Singular Value Decomposition (SVD) can be computationally prohibitive, scaling as $\mathcal{O}(L \cdot \min(D_1^2 D_2, D_1 D_2^2))$, or $\mathcal{O}(L \cdot D^3)$ for square weight matrices of hidden dimension $D$. To scale SR3-S to giant LLMs (e.g., $D = 4096$ or $D=8192$), practitioners can utilize randomized SVD or power iterations to approximate the spectral norm $\sigma_{\max}(V_k^{(l)})$ of each task vector offline. A power iteration algorithm applied to $V^T V$ converges rapidly to the largest singular value. Each step of power iteration requires only two matrix-vector multiplications, resulting in a complexity of $\mathcal{O}(m \cdot D^2)$ for $m$ iterations. For $m \le 3$, this reduces the computational footprint from $\mathcal{O}(D^3)$ to $\mathcal{O}(D^2)$, representing a massive speedup of $\mathcal{O}(D)$ (over three orders of magnitude for modern LLM hidden dimensions) and making spectral norm profiling exceptionally fast and scalable.
    \item \textbf{PAC-Bayesian Generalization Bounds:} While we focus on Rademacher complexity, a PAC-Bayesian framework could provide tighter, data-dependent generalization bounds for stochastic routing mechanisms (e.g., Gumbel-Softmax routers). Our initial theoretical discussion in Section 3.2 suggests that under stochastic ensembling, an asymmetric prior $P$ penalizing large task-vector norms can naturally derive the geometry-aware asymmetric scaling of SR3 directly from KL divergence minimization, providing a fertile ground for future stochastic ensembling theory.
    \item \textbf{Dynamic Projection Matrices:} We employ a frozen random projection matrix $P$ to map representations onto a unit sphere. Optimizing this projection matrix dynamically alongside the router parameters while maintaining a bounded lipschitz constant represents a promising direction for enhancing representation-space discriminability without sacrificing theoretical guarantees.
    \item \textbf{Hybrid and Adaptive Capacity Controllers:} To resolve the capacity-allocation trade-off (the over-repression of high-complexity expert task-specific accuracy, such as SVHN), future research should investigate adaptive capacity controllers that dynamically modulate the regularization force. For instance, the regularization multiplier for expert $k$ at step $t$ could be scaled inversely by the running average of the gradient norm for that expert's routing weights: $\lambda_k^{(t)} = \lambda_{\text{base}} \cdot v_k \cdot \exp\left(-\gamma \cdot \|\nabla_{W_k} \mathcal{L}_{\text{CE}}\|_2\right)$. This would allow the router to adaptively relax the strict learning-theoretic complexity bounds when calibration cues are exceptionally strong and unambiguous, successfully preserving the expert's specialization capacity on difficult task domains without sacrificing overall generalization boundaries.
\end{enumerate}

Ultimately, this work bridges the gap between empirical weight-space model merging and first-principles statistical learning theory. We hope our analysis inspires researchers to move away from heuristic ensembling and toward provably optimal, geometry-aware merging paradigms.
